\documentclass[14pt]{extarticle}
\usepackage[a4paper,margin=1in]{geometry}
\usepackage{times}
\usepackage{setspace}
\setstretch{1.05}
\pagestyle{empty}

\begin{document}

\begin{center}
\textbf{A List of Changes (Significant Improvement on Content)}\\
\textbf{Last updated date: 11 May 2026}
\end{center}

\vspace{0.5em}

\noindent
\textbf{The most significant change}

The revised paper substantially expands the evaluation metrics for the proposed
IPM method. The static and dynamic evaluations now report absolute Euclidean
error, axis-wise error, normalized percentage error, MAE, RMSE, and standard
deviation where appropriate, instead of relying only on average estimated
positions. This provides a more complete assessment of accuracy and error
variation under both stationary and walking conditions.

\vspace{0.8em}

\noindent
\textbf{The second most significant change}

The revised paper adds field-coordinate spatial analyses for both static and
dynamic evaluations. The static analysis compares reference object positions
with mean IPM estimates and error vectors for six field objects, while the
dynamic analysis shows filtered estimates over 10 walking trials with 15 cm
tolerance regions. These additions make the spatial behavior of the method and
the effect of walking motion clearer.

\vspace{0.8em}

\noindent
\textbf{The third most significant change}

The revised paper strengthens the comparison against the previously used
polynomial regression method. The comparison now includes distance-wise
estimated-distance and absolute-error trends, together with MAE, RMSE, and
maximum error over reference distances from 20 cm to 400 cm. The results show
that IPM reduces the tested MAE from 36.90 cm to 2.24 cm and avoids the
long-distance saturation observed in the regression method.

\vspace{1.2em}

\noindent
\textbf{Deletion if any}

No major technical content was deleted. The revision mainly adds evaluation
metrics, spatial evaluation analysis, and a stronger quantitative comparison.

\end{document}
